%LTeX: language=fr-FR
\documentclass{article}
\usepackage{tasks}
\usepackage{amsmath,amssymb,amsfonts}
\usepackage[french]{babel}
\usepackage{enumitem}
\usepackage{currfile}
\usepackage{tikz}
\usetikzlibrary{arrows.meta}
\usepackage{tcolorbox}
\usepackage{fancyhdr}
\usepackage{array}
\usepackage[headsep=3mm,footskip=5mm,includefoot,includehead]{geometry}

\setlength{\parindent}{0cm}

\pagestyle{empty}
\input{\currfiledir ../def.tex}
\def \Document {Activité 1}


\geometry{left=1.5cm, right=1.5cm, top=1cm, bottom=1cm}
\setlength{\headheight}{14.0pt}

\newcolumntype{C}[1]{>{\centering\arraybackslash}m{#1}}

% Droite graduée des tailles, de 0,9 m à 2,1 m
\newcommand{\droitetailles}{%
  \begin{tikzpicture}[x=11cm]
    \draw[-{Latex[length=2mm]}] (0.85,0) -- (2.15,0) node[right] {\small $x$ (m)};
    \foreach \x/\l in {0.9/0{,}9, 1.0/1{,}0, 1.1/1{,}1, 1.2/1{,}2, 1.3/1{,}3, 1.4/1{,}4, 1.5/1{,}5, 1.6/1{,}6, 1.7/1{,}7, 1.8/1{,}8, 1.9/1{,}9, 2.0/2{,}0, 2.1/2{,}1} {
      \draw (\x,0.1) -- (\x,-0.1);
      \node[below] at (\x,-0.1) {\scriptsize$\l$};
    }
    \foreach \x in {0.95,1.05,...,2.05} \draw (\x,0.05) -- (\x,-0.05);
  \end{tikzpicture}%
}

\begin{document}
\pagestyle{fancy}
\fancyhf{}
\fancyhead[L]{\textbf{\Niveau}}
\fancyhead[C]{\textbf{\ChapitreCourt}}
\fancyhead[R]{\textbf{\Document}}

\large

\begin{center}
  {\Large\bfseries Activité d'introduction : le parc d'attractions}
\end{center}

Dans un parc d'attractions, chaque manège impose des conditions de \textbf{taille}, notée $x$ (en mètres), pour des raisons de sécurité.

\begin{center}
  \renewcommand{\arraystretch}{1.5}
  \begin{tabular}{|C{3.2cm}|C{11cm}|}
    \hline
    \textbf{Manège} & \textbf{Condition d'accès} \\ \hline
    Le Grand Huit & mesurer au moins $1{,}40$ m \\ \hline
    Le Petit Train & mesurer strictement moins de $1{,}20$ m \\ \hline
    La Rivière & mesurer entre $1{,}10$ m et $1{,}95$ m (ces tailles comprises) \\ \hline
    La Chute Libre & mesurer strictement plus de $1{,}30$ m, et au plus $1{,}90$ m \\ \hline
  \end{tabular}
\end{center}

% =====================================================
\subsection*{Partie A : traduire les conditions}

\begin{enumerate}
  \item Traduire chaque condition d'accès par une ou deux inégalités portant sur $x$.

  \smallskip
  \begin{tasks}(2)
    \task Grand Huit : \dotfill
    \task Petit Train : \dotfill
    \task Rivière : \dotfill
    \task Chute Libre : \dotfill
  \end{tasks}

  \item Sur chacune des droites ci-dessous, colorier les tailles qui permettent d'accéder au manège. Indiquer par un symbole si les tailles aux extrémités sont autorisées ou non.

  \medskip
  \begin{tabular}{@{}p{2.3cm}l@{}}
    Grand Huit & \droitetailles \\[1.2em]
    Petit Train & \droitetailles \\[1.2em]
    Rivière & \droitetailles \\[1.2em]
    Chute Libre & \droitetailles
  \end{tabular}
\end{enumerate}

\begin{tcolorbox}[colback=white, colframe=black, title=Une nouvelle notation]
  L'ensemble de tous les nombres $x$ tels que $1{,}10 \leq x \leq 1{,}95$ se note $[1{,}10\,;\,1{,}95]$.

  L'ensemble de tous les nombres $x$ tels que $1{,}30 < x \leq 1{,}90$ se note $]1{,}30\,;\,1{,}90]$.

  Ces ensembles s'appellent des \textbf{intervalles}.
\end{tcolorbox}

\begin{enumerate}[resume]
  \item En observant ces deux exemples, expliquer ce qu'indique le sens des crochets.

  \medskip
  \dotfill

  \medskip
  \dotfill

  \item Pour le Grand Huit, il n'y a pas de taille maximale. On utilise alors le symbole $+\infty$ (\og plus l'infini \fg{}). Proposer une écriture de l'ensemble des tailles autorisées :
  \dotfill
  \item Proposer de même une écriture pour le Petit Train : \dotfill
\end{enumerate}

\newpage

% =====================================================
\subsection*{Partie B : qui peut monter ?}

Cinq amis se présentent à l'entrée du parc.

\begin{enumerate}[resume]
  \item Compléter le tableau par \og oui \fg{} ou \og non \fg{}.

  \medskip
  \begin{center}
    \renewcommand{\arraystretch}{1.6}
    \begin{tabular}{|C{3cm}|C{2.6cm}|C{2.6cm}|C{2.6cm}|C{2.6cm}|}
      \hline
      & \textbf{Grand Huit} & \textbf{Petit Train} & \textbf{Rivière} & \textbf{Chute Libre} \\ \hline
      Léo : $1{,}15$ m & & & & \\ \hline
      Inès : $1{,}40$ m & & & & \\ \hline
      Noé : $1{,}30$ m & & & & \\ \hline
      Sarah : $1{,}90$ m & & & & \\ \hline
      Yanis : $1{,}96$ m & & & & \\ \hline
    \end{tabular}
  \end{center}

  \item On dit que $1{,}40$ \textbf{appartient} à l'intervalle $[1{,}10\,;\,1{,}95]$ et on écrit $1{,}40 \in [1{,}10\,;\,1{,}95]$.
  Compléter avec $\in$ ou $\notin$ (\og n'appartient pas à \fg{}) :

  \begin{tasks}(3)
    \task $1{,}30 \ \dots\ ]1{,}30\,;\,1{,}90]$
    \task $1{,}90 \ \dots\ ]1{,}30\,;\,1{,}90]$
    \task $1{,}96 \ \dots\ [1{,}10\,;\,1{,}95]$
  \end{tasks}
\end{enumerate}

% =====================================================
\subsection*{Partie C : le pass \og Duo \fg{}}

\begin{enumerate}[resume]
  \item Le pass \og Duo \fg{} permet de faire la Rivière \textbf{et} le Grand Huit. Quelles tailles permettent de faire \textbf{les deux} manèges ? Répondre à l'aide d'un intervalle.

  \medskip
  \dotfill

  \item Une famille veut qu'au moins un manège soit accessible à chacun de ses enfants : le Petit Train \textbf{ou} le Grand Huit. Colorier sur la droite les tailles qui permettent de faire \textbf{au moins un} de ces deux manèges.

  \medskip
  \droitetailles

  \medskip
  Ces tailles forment-elles un seul intervalle ? \dotfill

  \item Existe-t-il une taille permettant de faire à la fois le Petit Train et le Grand Huit ? \dotfill
\end{enumerate}

\end{document}
